\section*{Team Contributions}

\begin{xltabular}{\linewidth}{@{}llX@{}}
  \caption{Team Contributions}
  \label{tab:contributions} \\
  \toprule
  Member & Student ID & Contribution \\
  \midrule
  \endfirsthead
  \multicolumn{3}{@{}l}{\normalfont Table \thetable. Team Contributions (continued)} \\
  \toprule
  Member & Student ID & Contribution \\
  \midrule
  \endhead
  \bottomrule
  \endlastfoot
  \linespread{1}\selectfont
  Aldo Kristian & 2532015 & Responsible for developing the Introduction and Conclusion sections. His work included explaining the research background, describing the research context and objectives, introducing the variables, and summarizing the main findings. He also helped connect the statistical results with the overall topic of gadget use as a learning medium. \\
  Jonatan & 2532035 & Responsible for organizing and directing the overall development of the article, including determining the structure, analysis flow, and consistency between sections. He prepared and organized the dataset, presented the raw data, and conducted the main statistical analysis using IBM SPSS. His work included calculating descriptive statistics, Pearson correlation, R\textsuperscript{2}, t-test, and F-test, establishing measurement criteria and benchmarks, creating graphs, and interpreting the statistical results and significance values. He also analyzed the relevance of the relationships between the variables, reviewed the cause-and-effect interpretation based on the research design, and ensured that the statistical findings were presented accurately throughout the article. \\
  Kevin Li & 2532042 & Responsible for developing the Data Visualization and Learning Reflection sections. His work included presenting the statistical findings through visualizations, explaining patterns observed in the scatter plot, and describing how visualization supports the interpretation of the relationship between the variables. He also prepared the Learning Reflection by summarizing the group's learning experience and explaining the importance of understanding and interpreting statistical results. \\
  Vincent Yongky Pratama & 2532039 & Responsible for developing the Basic Probability section and implementing the probability analysis using Python. His work included creating the Python code for data processing and probability calculations, defining the sample space and events, and calculating probability, intersection, and conditional probability. He also interpreted the probability results and connected the calculations to the 69 respondents to demonstrate how probability concepts can be applied to real research data. He also prepared the LaTeX version of the article, including its layout, tables, and bibliography formatting. \\
\end{xltabular}
